% =============================================================================
%  Capítulo 5 — Conclusiones (se escribe AL FINAL)
% =============================================================================
\chapter{Conclusiones}
\label{cap:conclusiones}

\todo{Estructura sugerida (3--5 páginas):

\paragraph{Resumen del trabajo realizado.}
Recapitular brevemente los objetivos del Cap.~\ref{cap:introduccion} y
relacionar cada uno con lo entregado.

\paragraph{Logros.}
\begin{itemize}
  \item Sistema modular \texttt{diffusion\_lib} que cubre los 5 puntos del
        enunciado.
  \item Implementación reproducible de VE/VP, EM/PC/PF-ODE, schedules
        lineal/coseno, métricas BPD/FID/IS, generación controlable.
  \item Validación cuantitativa con métricas estándar de la literatura.
\end{itemize}

\paragraph{Limitaciones.}
\begin{itemize}
  \item Datasets de baja resolución (MNIST, SVHN, gatos $32\times32$).
  \item U-Net sin las mejoras de las arquitecturas más recientes (DiT,
        EDM).
  \item Sin paralelismo de datos / multi-GPU.
  \item La calidad sobre imágenes naturales sigue lejos de los SOTA.
\end{itemize}

\paragraph{Trabajo futuro.}
\begin{itemize}
  \item Extender a resoluciones mayores ($64\times64$, $128\times128$)
        con arquitectura adaptada.
  \item Integradores de orden superior para PF-ODE (Heun, RK4, DPM-Solver).
  \item Métricas alternativas: CLIP-FID, Precision \& Recall.
  \item Modelos latentes (Latent Diffusion).
\end{itemize}

\paragraph{Reflexión personal.}
Lecciones aprendidas, dificultades técnicas encontradas, conexión con
otras asignaturas del grado.}
